\documentclass[10pt,a4paper]{amsart}
\usepackage[margin=19mm]{geometry}
\usepackage{amsmath,amssymb,mathtools}
\usepackage[scr=rsfs,cal=euler]{mathalfa}
\usepackage{xfrac}
\usepackage[unicode,hidelinks,bookmarks=false]{hyperref}
\newcommand{\ie}{i.e.\ }
\newcommand{\cf}{cf.\ }
\newcommand{\resp}{respectively\ }
\newcommand{\Subsection}{\subsection}
\providecommand{\nobreakdash}{\nobreak\hbox{-}}
\setlength{\emergencystretch}{2em}
\multlinegap0pt
\usepackage[shortlabels]{enumitem}
\usepackage{amssymb}
\usepackage{mathtools}
\newtheorem{theorem}{Theorem}
\newtheorem{lemma}{Lemma}
\newtheorem{proposition}{Proposition}
\usepackage{cleveref}
\crefname{theorem}{Theorem}{Theorems}
\crefname{lemma}{Lemma}{Lemmas}
\crefname{proposition}{Proposition}{Propositions}
\newcommand{\dist}{\operatorname{dist}}
\newcommand{\supp}{\operatorname{supp}}
\title{Repairing the refined-decoupling proof of the 5/4 planar pinned Falconer theorem}
\author{Shalender Singh, Vishnu Priya Singh Parmar}
\date{Source: arXiv:2608.28711v1 (CC BY 4.0)}
\begin{document}
\maketitle

\noindent\emph{Source and licence.} Repairing the refined-decoupling proof of the 5/4 planar pinned Falconer theorem. Version arXiv:2608.28711v1 (CC BY 4.0), distributed by arXiv under the Creative Commons Attribution 4.0 International licence, http://creativecommons.org/licenses/by/4.0/, as recorded in the arXiv licence metadata for this exact version and shown on the abstract page https://arxiv.org/abs/2608.28711v1. Full licence notice: \texttt{licenses/arxiv-2608.28711-CC-BY-4.0.txt}.\par\smallskip

\noindent\textit{Editorial scope.} Counted unit: the article's proposition prop:appendix-one-step (source0.tex lines 1493--1508). The article's complete original proof of it is delivered with it (source0.tex lines 1510--1657, one \texttt{proof} environment of 148 lines). No mathematics is written, repaired or re-derived: the statement and the proof are the article's own lines, in the article's own notation, and no retained line is substituted or re-flowed.  Retained uncounted as the source-local context the proof needs: lines 1231--1251; lines 1299--1304; lines 1309--1312; lines 1330--1333; lines 1351--1378; lines 1404--1406; lines 1423--1431.  Nothing was omitted from any retained span, so the package carries no omission record.  No retained line contains a citation, so the wrapper carries no bibliography and no bibliography entry.  No retained line contains a figure environment, an image-inclusion command or drawing code, so no image is delivered and the package declares no asset dependency.  Editorial formatting only, in addition: an AMS \texttt{amsart} preamble that restates the article's own declarations and macros used by the retained lines, namely the environments \texttt{theorem}, \texttt{lemma}, \texttt{proposition} on separate counters, together with the standard packages those lines need.  The retained environments print in this document as Theorem 1, Lemma 1, Proposition 1 in order of appearance, whereas the article prints the same items with the numbering of its own full document.  The complete native source of the article as distributed in the arXiv e-print for this exact version is bundled unchanged as \texttt{sources/208803/source0.tex}.
\begin{theorem}[Weighted $\ell^2$ decoupling input]
\label{thm:appendix-l2}
Fix $2\le p\le6$, $\epsilon'>0$, and $N_1\ge1$.  There is
$D=D(\epsilon',p,N_1,\{C_m\})$ such that the following holds for every phase
in the normalized rescaling family.  For $0<t\le1$, let $\mathcal J_t$ be a
bounded-overlap cover of $[-1,1]$ by intervals $J$ of length $2t^{1/2}$ and
put
\[
 \tau_J(t)=\{\xi:\xi_1\in J,\ |\xi_2-g(\xi_1)|\le t\}.
\]
If $F=\sum_{J\in\mathcal J_t}F_J$ and
$\supp\widehat F_J\subset\tau_J(t)$, then for every square $Q$ of side
$t^{-1}$,
\begin{equation}\label{eq:appendix-l2}
 \|F\|_{L^p(Q)}
 \le D t^{-\epsilon'}
 \left(\sum_{J\in\mathcal J_t}
       \|F_J\|_{L^p(w_Q)}^2\right)^{1/2},
\end{equation}
where $w_Q(x)=(1+t\dist(x,Q))^{-N_1}$.
\end{theorem}

\begin{equation}\label{eq:appendix-Adef}
 \left\|\sum_{T\in\mathcal W}F_T\right\|_{L^p(Y)}
 \le \mathfrak A(R,\Lambda,a,b)
 \left(\frac{M}{W}\right)^{\frac12-\frac1p}
 \left(\sum_{T\in\mathcal W}\|F_T\|_p^2\right)^{1/2}
\end{equation}

\begin{equation}\label{eq:appendix-multiplicity}
 \#\{T\in\mathcal W:bT\cap Q\ne\varnothing\}\le M
 \quad(Q\subset Y).
\end{equation}

\begin{equation}\label{eq:cap-in-tau}
 \theta_I^\Lambda
 \subset\tau_{J(I)}(R^{-1/2})
\end{equation}

\begin{lemma}[Exact parabolic rescaling]
\label{lem:appendix-rescaling}
Let $\widetilde g=g_{x_0,R}$.  Then:
\begin{enumerate}[label=\textup{(\alph*)},leftmargin=2em]
\item $\Phi(\Box_{J,u})=\mathcal B_{R'}$.  The inverse images of the
$R'^{1/2}$-squares are parallelograms of transverse width $R^{1/2}$ and
longitudinal length $R^{3/4}$.
\item If $I\subset J$ and $I'=R^{1/4}(I-x_0)$, then
\[
 \Psi(\theta_I^\Lambda)=(\theta_{I'}')^\Lambda,
 \qquad
 \Phi(sT_{I,v}\cap\Box_{J,u})=sT'_{I',v'},
 \quad v'=R^{-1/4}(v-u),
\]
where the primed objects are defined at scale $R'$ for $\widetilde g$ with
the same $\Lambda$.  The transformed intervals form another allowed translate of the
lower-scale lattice cover, with the same spacing parameter $\sigma_0$.
\item If
$H^\Phi(x')=(e^{-2\pi i(x_0,g(x_0))\cdot x}H(x))|_{x=\Phi^{-1}(x')}$,
then
\begin{align*}
 \supp\widehat{H^\Phi}&\subset\Psi(\supp\widehat H),\\
 \|H^\Phi\|_{L^p(\Phi(E))}
 &=|\det\Phi|^{1/p}\|H\|_{L^p(E)},
 \qquad |\det\Phi|=R^{-3/4}.
\end{align*}
\end{enumerate}
\end{lemma}

\begin{equation}\label{eq:aT-in-box}
 (a+1)T\subset\Box(T).
\end{equation}

\begin{lemma}[Transfer from a rescaling cell to the original square]
\label{lem:appendix-transfer}
Assume $R\ge2^{22}$, $64\le\Lambda\le R^{1/4}/100$, and
$C\cap Q''\ne\varnothing$ with $Q''\in\mathcal N(Q)$.  If
$b'T\cap C\ne\varnothing$, then the center of $Q$ lies in
\[
 (b'+12\Lambda^{-3/4})T.
\]
\end{lemma}

\begin{proposition}[One induction step]
\label{prop:appendix-one-step}
Let $R\ge2^{22}$,
$R^{\vartheta/2}\le\Lambda\le R^{1/4}/100$, $\Lambda\ge64$, and
\[
 1\le a\le10,
 \qquad a+12\Lambda^{-3/4}\le b\le20.
\]
Take $N_1=\lceil2000/\vartheta\rceil$ in
\cref{thm:appendix-l2}, and let $D$ be the corresponding constant.  Then
\begin{equation}\label{eq:appendix-recurrence}
 \mathfrak A(R,\Lambda,a,b)
 \le C D(\log R)^4R^{\epsilon'/2}\Lambda^{1/2}
 \mathfrak A(R^{1/2},\Lambda,a,b-12\Lambda^{-3/4})+1.
\end{equation}
\end{proposition}

\begin{proof}
Fix a configuration and abbreviate the right-hand side of
\eqref{eq:appendix-Adef}, without $\mathfrak A$, by $\mathrm{RHS}$.  All
errors below are chosen smaller than $R^{-40}\mathrm{RHS}$; the polynomial
number of such errors accounts for the final additive constant.

First pigeonhole the $R^{1/2}$-squares so that
$\|F\|_{L^p(Q)}$ is comparable on the retained union $Y_0$ and
\begin{equation}\label{eq:appendix-square-pigeonhole}
 \|F\|_{L^p(Y)}^p\lesssim(\log R)\|F\|_{L^p(Y_0)}^p.
\end{equation}
For a large interval $J$, let
$F_{\tau_J}=\sum_{\Box=\Box_{J,u}}F_\Box$.  By
\eqref{eq:cap-in-tau}, its Fourier support lies in
$\tau_J(R^{-1/2})$.  Apply \cref{thm:appendix-l2} with $t=R^{-1/2}$ on a
square $Q\subset Y_0$ and set
\[
 \mathcal N^*(Q)
 =\{x:\dist(x,Q)\le\Lambda^{1/4}R^{1/2}\}.
\]
Since $\Lambda\ge R^{\vartheta/2}$ and
$N_1\ge2000/\vartheta$, the decoupling weight satisfies
\[
 w_Q(x)\le \Lambda^{-N_1/4}\le R^{-250}
 \qquad(x\notin\mathcal N^*(Q)).
\]
Moreover, \eqref{eq:aT-in-box} and the packet-tail hypothesis imply
\[
 \|F_\Box\|_{L^p(\mathbb R^2\setminus\Box)}
 \le \sum_{T\in\mathcal W_\Box}
       \|F_T\|_{L^p(\mathbb R^2\setminus aT)}
 \le R^{-100}\sum_{T\in\mathcal W_\Box}\|F_T\|_p.
\]
Consequently a cylinder $\Box_{J,u}$ that does not meet
$\mathcal N^*(Q)$ contributes only a negligible weighted tail.

For fixed $J$, use the adapted coordinate
$y_1=x_1+g'(x_0)x_2$.  The $y_1$-diameter of
$\mathcal N^*(Q)$ is at most
\[
 3.1(2\Lambda^{1/4}+1)R^{1/2}\le R^{3/4}.
\]
The cylinders with label $J$ have $y_1$-width $50R^{3/4}$ and their
centres are spaced by $2R^{3/4}$.  Hence at most $27$ of them meet
$\mathcal N^*(Q)$.  The triangle inequality followed by Cauchy--Schwarz
therefore gives
\[
 \|F_{\tau_J}\|_{L^p(w_Q)}
 \le 6\left(\sum_{\Box=\Box_{J,u}}
          \|F_\Box\|_{L^p(w_Q)}^2\right)^{1/2}
 +\text{negligible}.
\]

Inside $\mathcal N^*(Q)$, decompose into the grid squares
$Q''\in\mathcal N(Q)$.  Outside that collection the weight-tail estimate
above applies, and hence
\[
 \|F_\Box\|_{L^p(w_Q)}^p
 \lesssim\sum_{Q''\in\mathcal N(Q)}
       \|F_\Box\|_{L^p(Q'')}^p+\text{negligible}.
\]
Each $Q''$ meets at most eight of the rescaling parallelograms
$C\subset\Box$, and these cells have disjoint interiors.  Thus
\[
 \|F_\Box\|_{L^p(Q'')}^p
 \le\sum_{C\cap Q''\ne\varnothing}
       \|F_\Box\|_{L^p(Q''\cap C)}^p.
\]
Insert the last three displays into the weighted decoupling inequality and
use, since $p\ge2$,
$(\sum_i a_i^p)^{1/p}\le(\sum_i a_i^2)^{1/2}$.  We obtain
\begin{equation}\label{eq:appendix-triples}
 \|F\|_{L^p(Q)}
 \lesssim 6D R^{\epsilon'/2}
 \left(
  \sum_{(\Box,Q'',C)}
  \|F_\Box\|_{L^p(Q''\cap C)}^2
 \right)^{1/2},
\end{equation}
up to negligible terms.  Here $Q''\in\mathcal N(Q)$,
$C\cap Q''\ne\varnothing$, and the number of possible $Q''$ for fixed $Q$
is $O(\Lambda^{1/2})$.  This is the fully expanded weight-tail and
cell-decomposition step used below.

Put $b'=b-12\Lambda^{-3/4}$.  A non-negligible term in
\eqref{eq:appendix-triples} has $m_{\Box,b'}(C)\ge1$.  Pigeonhole the triples
according to dyadic values
\[
 m_{\Box,b'}(C)\sim M',\qquad |\mathcal W_\Box|\sim W',
\]
and then pigeonhole the squares according to the number
$n(Q)$ of retained triples.  There are only $O((\log R)^3)$ choices.  We
obtain a union $Y_2\subset Y_0$ on which $n(Q)\sim n$ and
\begin{equation}\label{eq:appendix-pigeonholed}
 \|F\|_{L^p(Y)}^p
 \lesssim (\log R)^4
 (D R^{\epsilon'/2})^p n^{p/2-1}
 \sum_Q\sum_{(\Box,Q'',C)\in\mathcal T(Q)}
 \|F_\Box\|_{L^p(Q''\cap C)}^p.
\end{equation}
A fixed triple occurs for at most $O(\Lambda^{1/2})$ choices of $Q$.  For
fixed $(\Box,C)$ the pieces $Q''\cap C$ are disjoint.  Therefore
\begin{equation}\label{eq:appendix-summed-triples}
 \|F\|_{L^p(Y)}^p
 \lesssim (\log R)^4(DR^{\epsilon'/2})^p
 n^{p/2-1}\Lambda^{1/2}
 \sum_{\Box:\,|\mathcal W_\Box|\sim W'}
 \|F_\Box\|_{L^p(Y_{\Box,M'})}^p,
\end{equation}
where $Y_{\Box,M'}$ is the union of cells $C$ with
$m_{\Box,b'}(C)\sim M'$.

Apply \cref{lem:appendix-rescaling} to a fixed cylinder.  The rescaled
packets form a configuration at scale $R^{1/2}$ with the same $\Lambda$,
support dilation $a$, counting dilation $b'$, multiplicity $O(M')$, and
packet count $|\mathcal W_\Box|$.  Hence, writing
$\mathfrak A'=\mathfrak A(R^{1/2},\Lambda,a,b')$,
\[
 \|F_\Box\|_{L^p(Y_{\Box,M'})}
 \lesssim \mathfrak A'
 \left(\frac{M'}{|\mathcal W_\Box|}\right)^{\frac12-\frac1p}
 \left(\sum_{T\in\mathcal W_\Box}\|F_T\|_p^2\right)^{1/2}.
\]
The families $\mathcal W_\Box$ are disjoint.  Since their sizes are
comparable to $W'$, summing the preceding inequality to the $p$th power gives
\begin{equation}\label{eq:appendix-induction-sum}
 \sum_{\Box:\,|\mathcal W_\Box|\sim W'}
 \|F_\Box\|_{L^p(Y_{\Box,M'})}^p
 \lesssim
 (\mathfrak A')^p
 \left(\frac{M'}{W}\right)^{p/2-1}
 \left(\sum_{T\in\mathcal W}\|F_T\|_p^2\right)^{p/2}.
\end{equation}

It remains to transfer the local multiplicity.  A fixed cylinder occurs in
at most $O(\Lambda^{1/2})$ triples attached to one $Q$.  Thus at least
$cn\Lambda^{-1/2}$ distinct cylinders occur.  Choose one retained triple in
each.  By \cref{lem:appendix-transfer}, all $\sim M'$ tubes counted by that
triple have their $b$-dilate meeting $Q$.  The cylinder families are
disjoint, and \eqref{eq:appendix-multiplicity} gives
\begin{equation}\label{eq:appendix-nM}
 nM'\lesssim\Lambda^{1/2}M.
\end{equation}
Substitute \eqref{eq:appendix-induction-sum} and
\eqref{eq:appendix-nM} into \eqref{eq:appendix-summed-triples}, take the
$p$th root, and use $2\le p\le6$.  All powers of $\Lambda$ are bounded by
$\Lambda^{1/2}$, and the result is \eqref{eq:appendix-recurrence}.
\end{proof}

\end{document}
